% !TEX root
\documentclass[12pt, headinclude=true]{scrartcl}

\usepackage{lmodern}
\usepackage{babel}

\usepackage{graphicx}
\usepackage{enumerate}

\usepackage{amssymb}
\usepackage{amsmath}
\usepackage{amsthm}
\usepackage{dsfont}
\usepackage{bm}
\usepackage{algorithm}

\usepackage{subfig}
\usepackage{wrapfig}

\usepackage{todonotes}

\usepackage{hyperref}
\hypersetup{
  colorlinks   = true,
  urlcolor     = blue,
  linkcolor    = blue,
  citecolor   = blue
}
\usepackage[noabbrev, capitalize,nameinlink]{cleveref}
\usepackage[round, comma]{natbib}

% \input{latex_templates/layout.tex}
\input{latex_templates/microtype.tex}
\input{latex_templates/theorems.tex}

\input{latex_shortcuts/math.tex}
\input{latex_shortcuts/stats.tex}
\input{latex_shortcuts/tex.tex}

%------------------------------------------------------------------------------%

\title{Transformation Local-Likelihood Estimation of the Tail Copula Density}
\author{Thibault Vatter\\ University of Applied Sciences Western Switzerland}

\begin{document}

\maketitle

Let $U \sim C$ be a $d$-dimensional random vector with $C$ an absolutely
continuous copula $C$, $c$ its density.
The (lower) tail copula is the homogeneous limit function
\begin{align*}%\label{eq:tail_copula}
  \Lambda(s)
  =
  \lim_{q\rightarrow0}
  q^{-1}C(qs),
  \qquad s\in[0,\infty)^d,
\end{align*}
whenever the limit exists.
Assume that the restriction of $\Lambda$ to $(0,1)^d$ is absolutely
continuous with tail density
\begin{align*}
  \lambda(s)
  =
  \frac{\partial^d}{\partial s_1\cdots\partial s_d}\Lambda(s),
  \qquad s\in(0,1)^d.
\end{align*}
In what follows, we consider nonparametric estimators of $\lambda$ based on a probit transformation of the data, similar to \citet{geenens2017probit} for the copula density $c$.
The idea is to define a probit-transformed tail density
\begin{align*}
  g(x)
  =
  \lambda(\Phi(x))\phi(x),
  \qquad x\in\mathbb R^d,
\end{align*}
where $\Phi$ is the standard normal distribution function taken component-wise,
and $\phi(x)=\prod_{j=1}^d\phi(x_j)$ is the standard normal density.
The advantage of this transformation is that the domain of $g$ is $\mathbb R^d$, which allows to use standard kernel methods for unbounded domains.
The tail density $\lambda$ can be recovered from $g$ via the inverse transformation
\begin{align*}
  \lambda(s)
  =
  \frac{
    g(\Phi^{-1}(s))
  }{
    \phi(\Phi^{-1}(s))
  },
  \qquad s\in(0,1)^d.
\end{align*}
Let $\{U_i\}_{i=1}^n$ be i.i.d. observations from $C$ and $q_n\rightarrow0$ be a sequence of thresholds such that $nq_n\to\infty$.
For observations in the lower tail $Q_n= \{i:U_i\in[0,q_n]^d\}$, define their rescaled versions and probit transformations respectively as $V_i=U_i/q_n$ and $X_i=\Phi^{-1}(V_i)$,
where equations are meant component-wise.
An estimator of $g$ can be written as
\begin{align*}
  \hat g_0(x)
  =
  \frac{1}{nq_n}
  \sum_{i \in Q_n}
  K_h(x-X_i),
  \qquad x\in\mathbb R^d,
\end{align*}
with $ K_h(x)=h_n^{-d}K(x/h_n),$
$K:\mathbb R^d\to[0,\infty)$ a kernel satisfying
$\int K(x)\,dx=1$, and
where $h_n$ is a bandwidth such that $h_n\rightarrow0$ and $nq_nh_n^d\to\infty$.
The corresponding estimator of the lower tail density is then simply
\begin{align*}
  \hat\lambda_0(s)
  =
  \frac{
    \hat g_0(\Phi^{-1}(s))
  }{
    \phi(\Phi^{-1}(s))
  },
  \qquad s\in(0,1)^d.
\end{align*}

More generally, for an integer $p\ge0$, let $A_p = \{\alpha\in\mathbb N_0^d:|\alpha|\le p\}$ be the set of multi-indices of order at most $p$, and
$P_a(z)$ be the polynomial of degree at most $p$ in $d$ variables with coefficients $a=(a_\alpha)_{\alpha\in A_p}$, that is,
\begin{align*}
  P_a(z)
  =
  \sum_{\alpha\in A_p} a_\alpha z^\alpha,
  \qquad z\in\mathbb R^d.
\end{align*}
Then, at a fixed point $x\in\mathbb R^d$, approximate locally $\log g(y)
\approx
P_a(y-x)$ for $y$ in a neighborhood of $x$,
and define $\hat a_p(x)$ as a maximizer of
\begin{align*}
  \ell_x(a)
  =
  \sum_{i \in Q_n}
  K_h(X_i-x)P_a(X_i-x)
  -
  nq_n
  \int_{\mathbb R^d}
  K_h(y-x)\exp\{P_a(y-x)\}\,dy,
\end{align*}
over the set of coefficients for which the integral is finite.
The local log-polynomial estimator of the transformed tail density is
\begin{align*}
  \hat g_p(x)
  =
  \exp\{\hat a_{p,0}(x)\},
\end{align*}
where $\hat a_{p,0}(x)$ denotes the fitted intercept. The associated
estimator of the lower tail density is
\begin{align*}
  \hat\lambda_p(s)
  =
  \frac{
    \hat g_p(\Phi^{-1}(s))
  }{
    \phi(\Phi^{-1}(s))
  },
  \qquad s\in(0,1)^d.
\end{align*}
The case $p=0$, i.e. local constant, corresponds to the estimator described above.
The cases $p=1$ and $p=2$ give the local log-linear and local
log-quadratic estimators, respectively.

\bibliographystyle{apalike}
\bibliography{mybib}

\end{document}